% Einheit 2 von 4 – Sinus- und Kosinusfunktion und ihre Merkmale (bank/trigonometrische-funktionen/e2.jsonl, zusammenbau v0.1)
%% TODO Zweigzeile Teil 1: was hier gelernt wird (Satz in Schülersprache aus dem Einheitstitel) – Einheit 2
\einheitenkopf[e2]{Einheit 2 von 4 · Sinus- und Kosinusfunktion und ihre Merkmale}
\zweigzeile{ab Kl. 10 · keine P10-Aufgabe}
\setcounter{aufgabe}{15}

%% TODO Titel: Ich-kann-Satz fehlt in der Bank (Platzhalter: Kettenname) – Nr. 16
%% TODO Anweisung: ein Satz über der ersten Teilaufgabe fehlt in der Bank – Nr. 16
\begin{aufgabe}{Sinus- und Kosinusfunktion}
%% TODO \rechenplatz{n}: Schrittzahl des Grundfalls fehlt in der Bank (nur bei fester Schreibform, 2.3 b)
\begin{teile}
\teil Wo ist die Welle zu Ende? Welcher Abschnitt ist genau eine ganze Wiederholung?\\ \kreuz{von A bis B}\\ \kreuz{von A bis C}

\begin{ksys}[xmin=0,xmax=720,ymin=-1.5,ymax=1.5,xstep=90,ystep=0.5,ablesen]\funktion{sin(\x)}{f}\punkt{0}{0}{A}\punkt{180}{0}{B}\punkt{360}{0}{C}\end{ksys}
\teil Fülle die Wertetabelle für y = sin(x) aus (zwei Stellen).

\wertetabelle{x in °}{y}{0,45,90,135,180,225,270,315,360}
\teil Trage die Punkte der Tabelle ein und zeichne die Welle y = sin(x).

\wertetabelle[0,0.71,1,0.71,0,-0.71,-1,-0.71,0]{x in °}{y}{0,45,90,135,180,225,270,315,360} \begin{ksys}[xmin=0,xmax=360,ymin=-1.2,ymax=1.2,xstep=45,ystep=0.2]\end{ksys}
\teil Rechtsachse für einen ganzen Vollkreis ab 0°, 9 cm Platz – wie viele Grad je Zentimeter? Welche Zahlen schreibst du an? \leerfeld[°] je cm
\teil Lies sin 135° am Graphen ab (eine Stelle).

\begin{ksys}[xmin=0,xmax=360,ymin=-1.2,ymax=1.2,xstep=45,ystep=0.2,ablesen]\funktion{sin(\x)}{f}\end{ksys}
\leerfeld
\teil Lies ab: Bei welchen Winkeln im ersten Vollkreis ist sin x ≈ 0,7?

\begin{ksys}[xmin=0,xmax=360,ymin=-1.2,ymax=1.2,xstep=45,ystep=0.2,ablesen]\funktion{sin(\x)}{f}\end{ksys}
$x_1 \approx$ \leerfeld[°], $x_2 \approx$ \leerfeld[°]
\teil Liegt P(50|0,77) auf dem Graphen von y = sin(x)? Rechne im Grad-Modus. \janein
\teil Welche Werte kann y = sin(x) annehmen? \leerfeld
\teil Gib alle Nullstellen von y = sin(x) zwischen 0° und 720° an. Welche Regelmäßigkeit siehst du? \leerfeld
\teil Gib Hoch- und Tiefpunkt von y = sin(x) im ersten Vollkreis ab 0° an. H(\leerfeld | \leerfeld), T(\leerfeld | \leerfeld)
\teil Nach welchem Winkel wiederholt sich y = sin(x) zum ersten Mal ganz? \leerfeld[°]
\teil y = sin(x) zwischen 100° und 250° – steigt oder fällt die Kurve?\\ \kreuz{steigt}\kreuz{fällt}
\teil Beschreibe die Symmetrie der Sinuskurve und belege sie mit einem Wertepaar.
\teil Um wie viel Grad muss man die Sinuskurve nach links schieben, damit sie auf der Kosinuskurve liegt?

\begin{ksys}[xmin=0,xmax=360,ymin=-1.2,ymax=1.2,xstep=45,ystep=0.2,ablesen]\funktion{sin(\x)}{f}\funktion{cos(\x)}{g}\end{ksys}
\leerfeld[°]
\teil Welche Gleichung gehört zum Graphen?\\ \kreuz{y = sin(x)}\\ \kreuz{y = cos(x)}

\begin{ksys}[xmin=0,xmax=360,ymin=-1.2,ymax=1.2,xstep=45,ystep=0.2,ablesen]\funktion{cos(\x)}{h}\end{ksys}
\teil Fülle die Wertetabelle für y = sin(x) aus (zwei Stellen), zeichne den Graphen und gib Wertebereich, Nullstellen, Hoch- und Tiefpunkt, kleinste Periode und Symmetrie an.

\wertetabelle{x in °}{y}{-180,-135,-90,-45,0,45,90,135,180} \begin{ksys}[xmin=-180,xmax=180,ymin=-1.2,ymax=1.2,xstep=45,ystep=0.2]\end{ksys}
\end{teile}
\end{aufgabe}

%% TODO Titel: Ich-kann-Satz fehlt in der Bank (Platzhalter: Kettenname) – Nr. 17
%% TODO Anweisung: ein Satz über der ersten Teilaufgabe fehlt in der Bank – Nr. 17
\begin{aufgabe}{Sinus- und Kosinusfunktion – Fehler finden, Begründen, Darstellungswechsel}
\begin{teile}
\teil Mehmet zeichnet die Sinuskurve so, dass sie im Hochpunkt spitz zuläuft wie ein Dach. Finde den Fehler.
\teil Begründe, warum sich die Sinuskurve nach einem Vollkreis wiederholt.
\teil Gehört die Wertetabelle zu f oder zu g?\\ \kreuz{f}\kreuz{g}

\wertetabelle[1,0,-1,0,1]{x in °}{y}{0,90,180,270,360} \begin{ksys}[xmin=0,xmax=360,ymin=-1.2,ymax=1.2,xstep=45,ystep=0.2,ablesen]\funktion{sin(\x)}{f}\funktion{cos(\x)}{g}\end{ksys}
\end{teile}
\end{aufgabe}
